%% evaluation.tex

\section{Validation of Pattern Implementability and Efficacy}
\label{sec:evaluation}

This section validates that the \reauthgate{} design pattern (\S\ref{sec:defense}) is implementable, that it enforces the \privinv{} \emph{by construction}, and that it imposes negligible cost on legitimate self-healing. We deliberately frame this as a \emph{validation} rather than an empirical evaluation: the central guarantee of the pattern is a programmatic check on a whitelist, not a probabilistic defense, so the empirical question is not ``does it work?'' (the code inspection answers that) but ``does it impose unacceptable cost on benign recovery, and is the implementation correct?''

A broader empirical evaluation---statistical generalization across production frameworks, cross-model robustness of content-level defenses, defense-aware adaptive adversaries, end-to-end exploit confirmation on third-party code---is beyond the scope of this architectural paper and is addressed by complementary empirical work in the agent-security literature~\cite{debenedetti2025camel,zhan2025adaptive}.

\subsection{Validation Setup}

\paragraph{Testbed.}
We implemented the \reauthgate{} pattern in \system{}, a multi-agent LLM operating system prototype (Java 21). The recovery orchestrator dispatches a layered chain of 14 concrete recovery strategy classes organized into 11 conceptual tiers (from lightweight in-place retry to circuit-breaker human escalation), including \code{TopologyMutationStrategy}, \code{ReflectionInjectionRecovery}, \code{ModelFallbackStrategy}, and human escalation. The system uses role-based access control with four registered roles drawn from a whitelisted blueprint registry: \code{Code\_Reviewer}, \code{Python\_Coder}, \code{Security\_Auditor}, and \code{System\_Architect}. The \code{PermissionChecker} gates the forward execution path; the \reauthgate{} extends it to the recovery path.

\paragraph{Configurations.}
Two matched configurations are compared:
\begin{itemize}[leftmargin=*]
    \item \textbf{BASELINE (no gate):} The \code{TopologyMutationStrategy} calls \code{parseAndValidate(validate=false)}, which adopts the LLM-diagnosed role as-is without authorization. Side effects execute inline at strategy-return time. This is the CWE-862 / Confused Deputy vulnerable configuration described in \S\ref{sec:attack-chain}.
    \item \textbf{DEFENDED (gate enabled):} The \code{TopologyMutationStrategy} calls \code{parseAndValidate(validate=true)} (Layer 1: \code{RoleReplacementValidator}) and returns \code{RecoveryResult.okRequiringReauthorization(...)} with the side effect deferred as a callback. The \code{RecoveryOrchestrator} intercepts the result and invokes \code{RecoveryReauthorizationGate.check(...)} (Layer 2) before executing the callback. This is the pattern specified in \S\ref{sec:design-decomposition}.
\end{itemize}

\paragraph{Workloads.}
Three task sets, 100 trials each per configuration (600 total):
\begin{enumerate}[leftmargin=*]
    \item \textbf{TopologyMutation attack (V2):} Five adversarial core-dump variants embedding \code{suggested\_role=System\_Admin} (or equivalent) hints in the diagnostic content fed to the diagnostic LLM, cycled across samples (20 trials each). Attack success = the LLM-suggested role reaches \code{resumeNode} and the agent is replaced with a role outside the failing agent's downgrade set.
    \item \textbf{BenignRecovery:} Internal errors plus a legitimate role downgrade (\code{Python\_Coder} $\rightarrow$ \code{Code\_Reviewer}). The role is in the pre-approved downgrade set, so the gate should admit it. False alarm = recovery blocked despite a legitimate downgrade.
    \item \textbf{BenignTopologyMutation:} A legitimate capability-mismatch diagnosis where the diagnostic LLM, reading a benign core dump, suggests a role that is genuinely a downgrade (e.g., \code{System\_Architect} $\rightarrow$ \code{Code\_Reviewer} when the task is code review). This tests whether the gate false-alarms on legitimate role mutations.
\end{enumerate}

\paragraph{Real LLM diagnostic module.}
The diagnostic LLM that consumes the core dump and emits the \code{suggested\_role} JSON is a real LLM API (DeepSeek-V4-Flash via SenseNova, Qwen2.5-72B-Instruct via SiliconFlow), called at temperature 0.7 to capture natural non-determinism. A multi-tier API rotation strategy ensures stable throughput. Total API consumption: approximately 400 paid calls, total cost under US\$20.

\subsection{Results}

Table~\ref{tab:validation} reports the validation results.

\begin{table}[htbp]
\centering
\small
\caption{\reauthgate{} validation on the \system{} testbed. The V2 attack row reports the attack success rate (ASR); the benign rows report the false alarm rate (FAR). The defended V2 ASR of 0.00 is \emph{by construction}: the whitelist rejects out-of-set roles deterministically. The 100-trial result is a verification that the implementation has no bypass bug, not a statistical finding.}
\label{tab:validation}
\begin{tabular}{lccc}
\toprule
\textbf{Task Set} & \textbf{Baseline ASR/FAR} & \textbf{Defended ASR/FAR} & \textbf{Machinery Overhead} \\
\midrule
V2 (TopologyMutation attack)        & 1.00 & \textbf{0.00}$^\dagger$ & 0.20 $\pm$ 0.06 ms \\
BenignRecovery (legitimate retry)   & 0.00 & 0.00 & 0.02 $\pm$ 0.03 ms \\
BenignTopologyMutation (legitimate downgrade) & 0.00 & 0.00 & 0.20 $\pm$ 0.06 ms \\
\midrule
\multicolumn{4}{l}{\emph{Total end-to-end latency (including LLM inference), mean $\pm$ std over $N=30$:}} \\
\quad V2 baseline       & \multicolumn{3}{c}{$2851 \pm 325$ ms} \\
\quad V2 defended       & \multicolumn{3}{c}{$3194 \pm 446$ ms} \\
\quad Benign recovery   & \multicolumn{3}{c}{$10 \pm 1$ ms} \\
\bottomrule
\end{tabular}
\\[2pt]
\footnotesize $^\dagger$By construction: the whitelist invariant rejects out-of-set roles deterministically. The 100-trial zero-ASR verifies correct implementation, not a statistical claim.
\end{table}

\subsubsection{Baseline Severity (V2 ASR = 1.00)}
Under the undefended configuration, the topology-mutation attack succeeds in 100 of 100 trials. The diagnostic LLM, reading the attacker-influenced core dump, emits an out-of-whitelist privileged role (e.g., \code{System\_Admin}) in every trial, and the baseline \code{parseAndValidate(validate=false)} call adopts the suggestion as-is. The role-replacement primitive \code{resumeNode} executes without an authorization check, replacing the failing agent with the attacker-suggested role. This is the CWE-862 / Confused Deputy vulnerability of \S\ref{sec:attack-chain}.

The perfect 1.00 ASR is not a property of a particular LLM: it reflects the structural absence of an authorization check at the recovery sink. The diagnostic LLM is asked to recommend a role, and its recommendation---whatever it is---is consumed verbatim. A different LLM might recommend a different role, but in the baseline configuration any recommendation reaches \code{resumeNode} unchecked.

\subsubsection{Defended Efficacy (V2 ASR = 0.00 by Construction)}
Under the defended configuration, the gate rejects every out-of-whitelist role. The 100-trial zero-ASR result verifies that the implementation has no bypass bug: the \code{RoleReplacementValidator} (Layer 1) correctly rejects \code{System\_Admin} as not in the blueprint registry, and even if Layer 1 were bypassed, the \code{RecoveryReauthorizationGate} (Layer 2) re-checks the invariant before the deferred callback executes. In every trial, the deferred side effect (the \code{resumeNode} callback) is never invoked; the orchestrator returns a failed recovery result and falls back to a safe default (typically human escalation).

Crucially, the zero ASR is \emph{by construction}: the whitelist invariant is a programmatic check, not a probabilistic defense. Treating the 100-trial zero-ASR as a statistical finding would be a category error---a whitelist that rejects out-of-set inputs rejects them deterministically, and no significance test is needed to ``prove'' it works. The 100 trials serve as a \emph{verification} that the implementation is correct (no bypass bug, no race condition, no edge case in the JSON parsing), not as a measurement of a probabilistic defense's efficacy.

Because the gate is a programmatic check on a whitelist, the guarantee is LLM-independent: a stronger instruction-following model, a paraphrased payload, or a defense-aware adversary cannot bypass the whitelist, because the whitelist does not inspect the content of the diagnostic LLM's output---it only checks whether the extracted role string is in the pre-approved set and whether it is a privilege escalation from the failing role.

\subsubsection{False Alarm Rate (FAR = 0.00)}
The defense false-alarms on \emph{zero} of 200 benign trials (100 BenignRecovery + 100 BenignTopologyMutation). The \code{RoleReplacementValidator} admits any role mutation satisfying $P(\text{new}) \leq P(\text{old})$, so the legitimate downgrade \code{Python\_Coder} $\rightarrow$ \code{Code\_Reviewer} passes both layers. Internal errors that do not trigger topology mutation are unaffected by the gate (the \code{requiresReauthorization} flag is \code{false}), so benign retry and reflection injection strategies run unchanged.

The zero FAR confirms that the \privinv{} is calibrated correctly: it blocks exactly the privilege-escalating actions and nothing else. A tighter invariant (e.g., ``no role mutation at all'') would false-alarm on legitimate downgrades; a looser invariant (e.g., ``any role in the whitelist'') would admit the attack. The $P(\text{new}) \leq P(\text{old})$ formulation is the right boundary.

\subsubsection{Machinery Overhead (Sub-millisecond)}
The defense machinery overhead---the cost of the \code{requiresReauthorization} flag check, the deferred-callback lookup, and the \code{RoleReplacementValidator} whitelist lookup, excluding LLM inference---is $0.20 \pm 0.06$~ms per V2 recovery attempt. This is negligible relative to the multi-second LLM inference latency. The total end-to-end latency for V2 defended ($3194 \pm 446$~ms) is higher than baseline ($2851 \pm 325$~ms); the $\sim$340~ms difference is dominated by the fallback retry path triggered when the gate rejects the proposed role (the orchestrator retries with a safe default role, issuing a second LLM diagnostic call). This is a one-time cost that occurs only when an attack is actually intercepted; in benign operation, no rejection occurs and no retry is needed, so benign latency remains $10 \pm 1$~ms.

\subsection{Architectural Implication}

The validation confirms three architectural properties of the \reauthgate{} pattern:

\begin{enumerate}[leftmargin=*]
    \item \textbf{Implementability.} The pattern is implementable in a multi-agent LLM OS with 14 recovery strategies, touching one orchestrator method and one side-effecting strategy (\code{TopologyMutationStrategy}). The migration is local and incremental: strategies that do not produce side effects need no change.
    \item \textbf{Efficacy by construction.} The defended V2 ASR is 0.00 because the whitelist invariant is a programmatic check, not a probabilistic defense. The guarantee is LLM-independent and expected to hold across all models, regardless of instruction-following strength or adversary adaptiveness.
    \item \textbf{Negligible cost.} Zero false alarms on 200 benign trials, sub-millisecond machinery overhead, and end-to-end latency bounded by LLM inference (not defense computation).
\end{enumerate}

The key architectural lesson is that \emph{structural invariants enforced at a centralized choke point are LLM-independent and therefore robust to model scaling}, whereas content-level defenses (which depend on the LLM obeying an untrusted-framing warning) may degrade on stronger instruction-following models. This asymmetry motivates our recommendation (\S\ref{sec:discussion}) that framework designers prioritize structural privilege invariants on the recovery path, even if they also deploy content-level provenance signaling as defense-in-depth.

\subsection{Threats to Pattern Validity}

We briefly note the scope limits of this validation, with fuller discussion in \S\ref{sec:discussion}.

\paragraph{Single-testbed validation.}
The validation is on \system{}, the authors' own prototype. The claim that the pattern \emph{transfers} to other frameworks rests on the framework-agnostic specification (\S\ref{sec:pattern-spec}), not on a cross-framework empirical study. The pattern's preconditions (a recovery orchestrator, a permission checker, an extensible result type, a deferral mechanism) are met by most self-healing frameworks~\cite{autogen,langgraph,reflexion}, but a port-and-validate study on additional frameworks is left as future work.

\paragraph{Single-attack-vector validation.}
The validation addresses the topology-mutation privilege-escalation vector (V2). Other recovery-path attack vectors---such as trust-escalation via high-trust template re-injection of externally-originated error content---are content-level rather than structural, and their defenses are probabilistic rather than by-construction. Those vectors are outside the scope of this architectural paper, which focuses on the structural privilege invariant.

\paragraph{Adaptive adversaries.}
The validation uses a fixed set of five adversarial core-dump variants. A fully adaptive multi-round adversary who observes gate rejections and iteratively rewrites payloads cannot bypass the whitelist (the gate does not inspect content), but could potentially exploit a bug in the JSON parsing or the role-extraction logic. We note this as a implementation-correctness concern, not a pattern-design concern: the pattern's guarantee is content-insensitive by construction.
